\documentclass[11pt,letterpaper]{article}
% Shared layout. Compile each document from this directory.
\usepackage[margin=0.8in,headheight=14pt,headsep=16pt,footskip=22pt]{geometry}
\usepackage[T1]{fontenc}
\usepackage{lmodern,amsmath,amssymb,booktabs,enumitem,titlesec,fancyhdr,microtype,xcolor,hyperref}
\definecolor{inkblue}{RGB}{25,65,100}
\hypersetup{colorlinks=true,linkcolor=inkblue,urlcolor=inkblue,filecolor=inkblue,pdfauthor={AI-assisted course study notes}}
\setlength{\parindent}{0pt}
\setlength{\parskip}{5pt}
\setlength{\emergencystretch}{1.5em}
\setlength{\abovedisplayskip}{7pt}
\setlength{\belowdisplayskip}{7pt}
\titleformat{\section}{\normalsize\bfseries\color{inkblue}}{}{0pt}{}
\titlespacing*{\section}{0pt}{9pt}{3pt}
\setlist[itemize]{leftmargin=15pt,itemsep=2pt,topsep=3pt,parsep=0pt}
\setlist[enumerate]{leftmargin=17pt,itemsep=3pt,topsep=3pt,parsep=0pt}
\pagestyle{fancy}
\fancyhf{}
\fancyhead[L]{\small ECE 750T\enspace /\enspace Responsible AI}
\fancyhead[R]{\small Fall 2026\enspace /\enspace Study notes}
\fancyfoot[R]{\small\thepage}
\renewcommand{\headrulewidth}{0.3pt}
\newcommand{\lesson}[2]{%
  \hypersetup{pdftitle={ECE 750T: #1 -- #2}}%
  {\LARGE\bfseries #1}\par\vspace{3pt}
  {\large #2}\par\vspace{3pt}
  {\footnotesize AI-assisted reconstruction from course slides, not a lecture transcript.
  Discussion reasoning is inferred; classroom responses were not supplied.}\par
  \vspace{3pt}\hrule\vspace{5pt}}
\newcommand{\sources}[1]{\par\vspace{3pt}{\footnotesize\textbf{Slide trail.} #1
  Page numbers refer to PDF pages, not the slides' continuous printed numbering.}\par}
\newcommand{\Prob}{\mathbb{P}}
\newcommand{\E}{\mathbb{E}}
\newcommand{\Yhat}{\widehat{Y}}
\newcommand{\indep}{\mathrel{\perp\!\!\!\perp}}

\begin{document}
\lesson{Lecture 3}{Evaluation and auditing: what does a score establish?}

\textbf{The lesson.} An evaluation supports a claim only for its population,
conditions, access rules, and measurement procedure. Auditing adds an assessment
against expectations and a route to accountability. The lecture moves from
generalization to subgroup errors, distribution shift, contaminated benchmarks,
LLM judges, and the institutions that can act on findings.

\section*{Match the evaluation design to the deployment claim}
The No Free Lunch statement assumes all binary labelings
of a finite input space are equally likely; expected error on unseen inputs is
then 50\%. This does not imply that structured real tasks are unlearnable.

An i.i.d. test estimates performance when train and test draws share a distribution.
Deployment may violate that premise. The diagrams distinguish changes in input
distribution (\emph{covariate shift}, with $P(Y\mid X)$ unchanged), corrupted
labels, changes in the feature--outcome relationship (\emph{concept shift}), and
changes in subgroup proportions. Different shifts call for different evidence.

The WILDS examples make the unit of generalization concrete: new camera locations,
hospitals, times, or regions. Training at hospitals A--C, validating at D, and
testing at E tests transfer to a held-out site. Randomly distributing images from
all hospitals across splits mainly tests new images from already represented sites.
Neither split answers every deployment question; choose according to the claim.

Reweighting changes the influence of observed cases and can target a different
mixture when suitable support and assumptions hold. It cannot supply missing kinds
of examples. The deep-learning illustration also cautions that changing weights
need not change the final decision boundary as expected throughout training.

\section*{Gender Shades: average performance can hide the central failure}
The external study built the Pilot Parliaments Benchmark: 1,270 public images from
six countries, evaluated by skin-type and perceived-gender groups. In the original
results, darker-female subgroup errors were about 21--35\%, while lighter-male
errors were at most 0.8\%. The intersection reveals a disparity that an overall
average, or a single coarse breakdown, can obscure.

The follow-up reported substantial improvements, with some disparities remaining.
This shows why external evaluation and public reporting can matter; it does not
establish that every audit causes improvement or that better classification
settles whether classifying gender from faces is an appropriate task. Keep the
study's labels, population, and normative question separate.

\section*{Prevent information leakage at the relevant boundary}
Fit preprocessing on training data; use validation for development decisions;
use the test set after those decisions are fixed. Separate related records by
person, site, or time when needed for the claim. Patient overlap, post-outcome
features, and test-informed preprocessing admit information unavailable in the
intended prediction setting. A nominally separate test file is insufficient.

\newpage
\section*{A benchmark can be compromised through several routes}
\textbf{Training exposure:} benchmark examples enter pretraining, undermining
held-out measurement. \textbf{Selection exposure:} repeated model or prompt
selection exploits benchmark feedback. \textbf{Evaluation-time access:} a tool-using
agent retrieves answers or reference solutions. These are different mechanisms;
an evaluation must specify the information and actions allowed at each stage.
Web-scale training and empirical scaling curves do not resolve provenance.

The slides' answer-seeking incident illustrates the last route: a task score
can reflect information access as well as the intended capability. Agent evaluation
must therefore examine behaviour and access boundaries alongside answer correctness.

\section*{Open-ended outputs require an explicit judgement procedure}
For summaries or dialogue there may be many acceptable outputs. Specify factual
accuracy, coverage, relevance, clarity, or harms; identify whose judgement is being
represented and for which use. A domain expert, user, annotator, and LLM judge
need not measure the same thing.

LLM judging provides scale, but needs validation. In the position-bias example,
swapping two fixed answers changes which wins. Order balancing or repeated
comparisons can reveal this sensitivity; agreement alone does not validate criteria.

Chatbot Arena measures preferences under its prompts, users, and comparison
process. The ``Leaderboard Illusion'' discussion raises selection and opportunity
concerns: more submitted variants and comparisons can increase the best observed
score. The slides also give Arena's response disputing preferential access.
This methodological dispute does not establish that the association is causal.

\section*{When does evaluation become an audit?}
Evaluation measures behaviour under specified conditions. An audit systematically
assesses a specified system or process against stated expectations, with sufficient
independence to support accountability. Subgroup error measurement can form part
of an audit; it needs scope, criteria, evidence, reporting, and responsible actors
to support a broader conclusion and corrective action.

Internal auditors can access development records and intervene early, but need
freedom to challenge developers. External auditors can question the organization's
framing, but may lack access. \textbf{Access, independence, and authority are
separate properties.} Funding, publication rights, conflicts of interest,
technical competence, affected-community participation, and powers to require
change determine whether findings can matter. Embedded evaluators gain access;
that alone does not establish independence or enforcement power.

A response to ``who should audit?'' should allocate roles: affected people identify
problems; evaluators design tests; auditors assess evidence; decision makers require
changes; oversight bodies review compliance. Specify who chooses the standard and
can publish an unwelcome finding.

\textbf{Midterm emphasis (inference).} Practise diagnosing leakage and shift,
designing the right split, explaining an intersectional audit result, and identifying
what evidence and authority an audit lacks. Judge bias and benchmark exposure are
strong conceptual candidates. Recalling leaderboard ranks or news dates has less
support as a study priority than explaining the failure mechanism.

\sources{\href{../L03.pdf}{L03.pdf}, pp.~3--8 (generalization/Gender Shades),
9--17 (shift and hygiene), 18--25 (benchmark exposure and judges),
26--32 (auditing); annotated copy consulted.}
\end{document}
